\documentclass[11pt]{article}

\usepackage[utf8]{inputenc}
\usepackage[T1]{fontenc}
\usepackage{lmodern}
\usepackage{geometry}
\usepackage{amsmath,amssymb,amsfonts,amsthm,mathtools}
\usepackage{bm}
\usepackage{enumitem}
\usepackage{microtype}
\usepackage{hyperref}
\usepackage{titlesec}
\usepackage{booktabs}
\usepackage{array}
\usepackage{setspace}
\usepackage{mathrsfs}
\usepackage{physics}

\geometry{margin=1in}
\hypersetup{colorlinks=true, linkcolor=black, urlcolor=blue, citecolor=black}
\setlist[enumerate]{topsep=0.4em,itemsep=0.35em}
\setlist[itemize]{topsep=0.3em,itemsep=0.25em}
\titleformat{\section}{\large\bfseries}{\thesection.}{0.5em}{}
\titleformat{\subsection}{\normalsize\bfseries}{\thesubsection.}{0.5em}{}
\setstretch{1.06}

\newcommand{\Aether}{\AE{}ther}
\newcommand{\Aflow}{\AE{}ther-flow}
\newcommand{\TheoryName}{\Aflow{} Interpretation of Relativity}
\newcommand{\TheoryProgram}{\Aether{} / \Aflow{} framework}

\newcommand{\CompletedMetric}{g^{\sharp}}

\newcommand{\Car}{\mathrm{car}}
\newcommand{\Cur}{\mathrm{cur}}
\newcommand{\Hom}{\mathrm{hom}}

\newcommand{\ForSpace}[1]{\mathcal I_{#1}^{\mathrm{for}}}
\newcommand{\PrimShell}{\mathcal S_{\mathrm{prim}}}
\newcommand{\ObsShell}{\mathcal S_{\mathrm{obs}}}

\newcommand{\Reduction}[1]{\mathcal R_{#1}}
\newcommand{\ReductionTriple}{\mathcal R_{\mathrm{tot}}}

\newcommand{\CurrentCarrier}[1]{\mathfrak C_{#1}^{\mathrm{cur}}}
\newcommand{\EqCarrier}[1]{\mathfrak C_{#1}^{\mathrm{eq}}}
\newcommand{\CurrentExtract}[1]{\mathscr E_{#1}^{\mathrm{cur}}}
\newcommand{\EqExtract}[1]{\mathscr E_{#1}^{\mathrm{eq}}}
\newcommand{\PrimCurrent}[1]{\mathcal J_{#1}^{\mathrm{prim}}}
\newcommand{\PrimTheta}[1]{\Theta_{#1}^{\mathrm{prim}}}
\newcommand{\EqProj}[1]{P_{#1}^{\mathrm{eq}}}

\newcommand{\DefectSpace}{\mathcal D}
\newcommand{\PrimDefectTuple}{\mathbf d^{\mathrm{prim}}}
\newcommand{\DeepDefectTuple}{\mathbf d^{\mathrm{deep}}}

\newcommand{\PrimResidualSector}{\mathfrak R_{\mathrm{res}}^{\mathrm{prim}}}
\newcommand{\PrimCompletionSector}{\mathfrak R_{\mathrm{cmp}}^{\mathrm{prim}}}
\newcommand{\PrimMatterSector}{\mathfrak R_{\mathrm{mat}}^{\mathrm{prim}}}
\newcommand{\DeepResidualSector}{\mathfrak R_{\mathrm{res}}^{\mathrm{deep}}}
\newcommand{\DeepCompletionSector}{\mathfrak R_{\mathrm{cmp}}^{\mathrm{deep}}}
\newcommand{\DeepMatterSector}{\mathfrak R_{\mathrm{mat}}^{\mathrm{deep}}}
\newcommand{\ResidualRead}{\mathscr R_{\mathrm{res}}}
\newcommand{\CompletionRead}{\mathscr R_{\mathrm{cmp}}}
\newcommand{\MatterRead}{\mathscr R_{\mathrm{mat}}}

\newcommand{\MetricOut}{\mathscr G_{\mathrm{out}}}
\newcommand{\MatterOut}{\mathscr T_{\mathrm{out}}}
\newcommand{\DeepMetricOut}{\mathscr G_{\mathrm{deep}}}
\newcommand{\DeepMatterOut}{\mathscr T_{\mathrm{deep}}}

\newcommand{\DeepSpace}[1]{\mathcal U_{#1}}
\newcommand{\DeepShell}{\mathcal U_{\mathrm{tot}}}
\newcommand{\VisibleMap}[1]{\mathcal V_{#1}}
\newcommand{\DeepCurrentCarrier}[1]{\mathfrak C_{#1}^{\mathrm{deep,cur}}}
\newcommand{\DeepEqCarrier}[1]{\mathfrak C_{#1}^{\mathrm{deep,eq}}}
\newcommand{\ChainSelect}[1]{\mathcal S_{#1}}
\newcommand{\ChainSelectTriple}{\mathcal S_{\mathrm{tot}}}

\newtheorem{definition}{Definition}
\newtheorem{lemma}{Lemma}
\newtheorem{proposition}{Proposition}
\newtheorem{remark}{Remark}
\newtheorem{corollary}{Corollary}
\newtheorem{theorem}{Theorem}

\title{\textbf{The \TheoryName{}}\\[0.4em]
\large Primitive-Reservoir Observer-Reduction Chain\\
\large Forcing Theorem}
\author{}
\date{}

\begin{document}

\maketitle

\begin{abstract}
After the primitive-reservoir observer-sector defect-fiber factorization theorem and the primitive-reservoir one-metric / ordinary-matter output map theorem, the remaining front-line burden on the active one-metric line is no longer another datum-level paper on the same primitive observer-reduction layer. The open question is why that now-recorded primitive observer-reduction chain exists at all.

The present note records a disciplined forcing theorem of exactly that kind. It introduces a chain-faithful deeper substrate package consisting of sectorwise deeper state spaces, visible maps, deeper current and equilibrium carrier maps, and forcing selectors into the recorded primitive forcing shells. Under a separation hypothesis that primitive shell points are already uniquely determined by their visible data and shell carriers, those selectors and hence the total primitive observer-reduction chain are uniquely induced by deeper explicit substrate structure. The current / projected-equilibrium package, the defect-fiber factorization of observer-sector outputs, and the metric / matter output maps all become pullbacks of that deeper package.

The benchmark claim boundary remains fixed. The deeper forcing still yields exactly one operative metric \(\CompletedMetric\) and exactly the same ordinary matter route through that metric; it introduces neither a second operative metric nor an enlarged ordinary matter law. This is not yet a completed first-principles derivation of GR from final substrate microphysics. Its narrower positive result is that the present primitive observer-reduction chain is justified as a forced output of deeper explicit substrate structure rather than merely recorded as fixed datum.
\end{abstract}

\section{Introduction}

The active derivational program of \TheoryName{} has now passed a natural datum-level stopping point. The primitive-reservoir observer-reduction datum origin theorem removed the current / equilibrium part of the primitive observer-reduction datum from the list of still-primitive blocks. The later defect-fiber factorization theorem showed that the observer-sector outputs already descend uniquely through the primitive defect fibers. The later output-map theorem then settled the remaining operative metric and ordinary-matter maps on that same chain, while keeping the benchmark one-metric package fixed.

That changes the routing logic. Once the current primitive observer-reduction chain has been recorded in full, another note at the same layer that merely relays the same one operative metric and the same ordinary matter route is no longer the front-line task. The remaining honest burden is deeper: explain why that entire chain itself is forced by explicit substrate structure.

The present note addresses exactly that burden and no more. It does not claim a final microphysical derivation of Einstein gravity. It does not enlarge the low-energy matter law, and it does not reopen the benchmark exact-closure package. Its claim is narrower and more disciplined. It shows that if one specifies a chain-faithful deeper substrate package whose visible maps and carrier data induce the recorded primitive shell states uniquely, then the current primitive observer-reduction chain is no longer separately chosen. It is forced.

\paragraph{Framework and claim boundary.}
\TheoryProgram{} denotes the broader \Aether{} / \Aflow{} framework built on the claims that the \Aether{} is the underlying four-dimensional substrate of reality, the \Aflow{} is its intrinsic ordered motion, observed three-dimensional space is the local experiential slice of that deeper substrate, S-time is the experienced order of change arising from matter, light, and the \Aflow{}, and observed expansion is the three-dimensional appearance of deeper four-dimensional ordered motion. Within that framework, \TheoryName{} denotes the exact-closure theory in which the effective gravitational dynamics are adopted to be exactly Einsteinian with universal matter coupling. In the active sequence, \emph{adoption} means use of that established relativistic dynamics without claiming substrate derivation, while \emph{derivation} is reserved for a first-principles recovery from explicit substrate variables.

The present note stays strictly inside that boundary. It does not claim that final substrate microphysics has already been found. Its narrower theorem is only that the now-recorded primitive observer-reduction chain can be forced from deeper explicit substrate structure while preserving the same benchmark one-metric package without change.

\section{Fixed Continuation Data After the Output-Map Theorem}

\begin{definition}[Recorded primitive observer-reduction chain]
\label{def:fixed}
Fix the following continuation data from the active primitive-reservoir line.
\begin{enumerate}
    \item The sectorwise forcing shells
    \[
    \ForSpace{\Car},\qquad
    \ForSpace{\Cur},\qquad
    \ForSpace{\Hom},
    \]
    their product
    \[
    \PrimShell
    \equiv
    \ForSpace{\Car}\times\ForSpace{\Cur}\times\ForSpace{\Hom},
    \]
    the observer-shell product
    \[
    \ObsShell
    \equiv
    X_{\Car}\times X_{\Cur}\times X_{\Hom},
    \]
    and the recorded shell reduction maps
    \[
    \Reduction{\Car}:\ForSpace{\Car}\to X_{\Car},\qquad
    \Reduction{\Cur}:\ForSpace{\Cur}\to X_{\Cur},\qquad
    \Reduction{\Hom}:\ForSpace{\Hom}\to X_{\Hom}.
    \]
    Define the total shell reduction by
    \[
    \ReductionTriple(w_{\Car},w_{\Cur},w_{\Hom})
    \equiv
    \bigl(
    \Reduction{\Car}(w_{\Car}),
    \Reduction{\Cur}(w_{\Cur}),
    \Reduction{\Hom}(w_{\Hom})
    \bigr).
    \]
    \item The already derived shell carriers and fixed shell extractors
    \[
    \CurrentCarrier{\sigma}:\ForSpace{\sigma}\to Z_{\sigma}^{\mathrm{cur}},
    \qquad
    \EqCarrier{\sigma}:\ForSpace{\sigma}\to Z_{\sigma}^{\mathrm{eq}},
    \]
    \[
    \CurrentExtract{\sigma}:Z_{\sigma}^{\mathrm{cur}}\to Y_{\sigma}^{\mathrm{cur}},
    \qquad
    \EqExtract{\sigma}:Z_{\sigma}^{\mathrm{eq}}\to Y_{\sigma}^{\mathrm{eq}},
    \qquad
    \sigma\in\{\Car,\Cur,\Hom\},
    \]
    together with the recorded shell-level current and equilibrium maps
    \begin{equation}
    \PrimCurrent{\sigma}
    =
    \CurrentExtract{\sigma}\circ\CurrentCarrier{\sigma},
    \qquad
    \PrimTheta{\sigma}
    =
    \EqExtract{\sigma}\circ\EqCarrier{\sigma}.
    \label{eq:carriercomposites}
    \end{equation}
    Also fix the compatibility projectors
    \[
    \EqProj{\Car},\qquad
    \EqProj{\Cur},\qquad
    \EqProj{\Hom}.
    \]
    \item The primitive defect tuple
    \[
    \PrimDefectTuple:\PrimShell\to\DefectSpace
    \]
    defined by
    \begin{equation}
    \begin{aligned}
    \PrimDefectTuple(w_{\Car},w_{\Cur},w_{\Hom})
    \equiv\;
    \Bigl(
    &\PrimCurrent{\Car}(w_{\Car}),
    \PrimCurrent{\Cur}(w_{\Cur}),
    \PrimCurrent{\Hom}(w_{\Hom}),
    \\
    &(\mathbf 1-\EqProj{\Car})\PrimTheta{\Car}(w_{\Car}),
    (\mathbf 1-\EqProj{\Cur})\PrimTheta{\Cur}(w_{\Cur}),
    \\
    &
    (\mathbf 1-\EqProj{\Hom})\PrimTheta{\Hom}(w_{\Hom})
    \Bigr).
    \end{aligned}
    \label{eq:primtuple}
    \end{equation}
    \item The already recorded primitive observer-sector outputs
    \[
    \PrimResidualSector:\PrimShell\to\mathcal V_{\mathrm{res}},
    \qquad
    \PrimCompletionSector:\PrimShell\to\mathcal V_{\mathrm{cmp}},
    \qquad
    \PrimMatterSector:\PrimShell\to\mathcal V_{\mathrm{mat}},
    \]
    together with the unique factor maps
    \[
    \ResidualRead,\qquad
    \CompletionRead,\qquad
    \MatterRead
    \]
    on \(\operatorname{im}\PrimDefectTuple\) such that
    \begin{equation}
    \PrimResidualSector=\ResidualRead\circ\PrimDefectTuple,\qquad
    \PrimCompletionSector=\CompletionRead\circ\PrimDefectTuple,\qquad
    \PrimMatterSector=\MatterRead\circ\PrimDefectTuple.
    \label{eq:sectorfactor}
    \end{equation}
    \item The recorded metric and ordinary matter outputs
    \[
    \MetricOut:\PrimShell\to\{\text{observer metrics}\},
    \qquad
    \MatterOut:\PrimShell\times\Psi\to\{\text{observer stress tensors}\},
    \]
    satisfying
    \begin{equation}
    \MetricOut(\mathbf w)=\CompletedMetric,
    \qquad
    \MatterOut(\mathbf w,\psi)
    =
    T_{\mu\nu}^{\mathrm{matter}}[\CompletedMetric,\psi]
    \label{eq:fixedoutputs}
    \end{equation}
    for every \(\mathbf w\in\PrimShell\) and \(\psi\in\Psi\). In particular, the recorded primitive observer-reduction chain introduces no second operative metric and no enlarged ordinary matter law.
\end{enumerate}
\end{definition}

\begin{remark}
Definition \ref{def:fixed} is the datum-level stopping point reached by the current primitive observer-reduction line. The purpose of the present note is not to restate that package as an acceptable datum. It is to ask whether the same package is itself induced by a deeper explicit substrate layer.
\end{remark}

\section{Chain-Faithful Deeper Substrate Forcing}

\begin{definition}[Chain-faithful deeper substrate package]
\label{def:deep}
For each sector label \(\sigma\in\{\Car,\Cur,\Hom\}\), a \emph{chain-faithful deeper substrate package} consists of:
\begin{enumerate}
    \item a deeper substrate state space
    \[
    \DeepSpace{\sigma};
    \]
    \item a visible map
    \[
    \VisibleMap{\sigma}:\DeepSpace{\sigma}\to X_{\sigma};
    \]
    \item deeper current and equilibrium carriers
    \[
    \DeepCurrentCarrier{\sigma}:\DeepSpace{\sigma}\to Z_{\sigma}^{\mathrm{cur}},
    \qquad
    \DeepEqCarrier{\sigma}:\DeepSpace{\sigma}\to Z_{\sigma}^{\mathrm{eq}};
    \]
    \item a forcing selector
    \[
    \ChainSelect{\sigma}:\DeepSpace{\sigma}\to\ForSpace{\sigma}
    \]
    such that
    \begin{equation}
    \Reduction{\sigma}\circ\ChainSelect{\sigma}
    =
    \VisibleMap{\sigma},
    \label{eq:visiblecompat}
    \end{equation}
    \begin{equation}
    \CurrentCarrier{\sigma}\circ\ChainSelect{\sigma}
    =
    \DeepCurrentCarrier{\sigma},
    \qquad
    \EqCarrier{\sigma}\circ\ChainSelect{\sigma}
    =
    \DeepEqCarrier{\sigma};
    \label{eq:carriercompat}
    \end{equation}
    \item the chain-separation property: if \(w,w'\in\ForSpace{\sigma}\) satisfy
    \begin{equation}
    \Reduction{\sigma}(w)=\Reduction{\sigma}(w'),
    \qquad
    \CurrentCarrier{\sigma}(w)=\CurrentCarrier{\sigma}(w'),
    \qquad
    \EqCarrier{\sigma}(w)=\EqCarrier{\sigma}(w'),
    \label{eq:separationdata}
    \end{equation}
    then \(w=w'\).
\end{enumerate}
Define the total deeper shell by
\[
\DeepShell
\equiv
\DeepSpace{\Car}\times\DeepSpace{\Cur}\times\DeepSpace{\Hom},
\]
and define the total forcing selector by
\[
\ChainSelectTriple(u_{\Car},u_{\Cur},u_{\Hom})
\equiv
\bigl(
\ChainSelect{\Car}(u_{\Car}),
\ChainSelect{\Cur}(u_{\Cur}),
\ChainSelect{\Hom}(u_{\Hom})
\bigr).
\]
\end{definition}

\begin{remark}
The positive move in Definition \ref{def:deep} is not merely to rename the primitive shell. The point is that visible data and explicit deeper carrier data on \(\DeepSpace{\sigma}\) induce the primitive shell state through a selector that is answerable to the recorded chain and can be proved unique under \eqref{eq:separationdata}.
\end{remark}

\begin{proposition}[Sectorwise forcing selectors are unique]
\label{prop:selector_unique}
Fix \(\sigma\in\{\Car,\Cur,\Hom\}\), and assume Definition \ref{def:deep}. If
\[
\widetilde{\mathcal S}_{\sigma}:\DeepSpace{\sigma}\to\ForSpace{\sigma}
\]
is another map satisfying
\[
\Reduction{\sigma}\circ\widetilde{\mathcal S}_{\sigma}
=
\VisibleMap{\sigma},
\qquad
\CurrentCarrier{\sigma}\circ\widetilde{\mathcal S}_{\sigma}
=
\DeepCurrentCarrier{\sigma},
\qquad
\EqCarrier{\sigma}\circ\widetilde{\mathcal S}_{\sigma}
=
\DeepEqCarrier{\sigma},
\]
then
\[
\widetilde{\mathcal S}_{\sigma}=\ChainSelect{\sigma}.
\]
\end{proposition}

\begin{proof}
Let \(u\in\DeepSpace{\sigma}\). By the displayed compatibility identities,
\[
\Reduction{\sigma}(\widetilde{\mathcal S}_{\sigma}(u))
=
\VisibleMap{\sigma}(u)
=
\Reduction{\sigma}(\ChainSelect{\sigma}(u)),
\]
\[
\CurrentCarrier{\sigma}(\widetilde{\mathcal S}_{\sigma}(u))
=
\DeepCurrentCarrier{\sigma}(u)
=
\CurrentCarrier{\sigma}(\ChainSelect{\sigma}(u)),
\]
and
\[
\EqCarrier{\sigma}(\widetilde{\mathcal S}_{\sigma}(u))
=
\DeepEqCarrier{\sigma}(u)
=
\EqCarrier{\sigma}(\ChainSelect{\sigma}(u)).
\]
The chain-separation property \eqref{eq:separationdata} therefore yields
\[
\widetilde{\mathcal S}_{\sigma}(u)=\ChainSelect{\sigma}(u).
\]
Since \(u\) was arbitrary, the two selectors agree everywhere.
\end{proof}

\begin{corollary}[The total forcing selector is unique]
\label{cor:selector_total}
For a chain-faithful deeper substrate package, the total selector \(\ChainSelectTriple:\DeepShell\to\PrimShell\) is uniquely determined by the visible maps and deeper carriers.
\end{corollary}

\begin{proof}
Apply Proposition \ref{prop:selector_unique} sectorwise and take the product.
\end{proof}

\section{Induced Current / Equilibrium Datum and Deeper Defect Tuple}

\begin{proposition}[The current / equilibrium part of the chain is induced from deeper carriers]
\label{prop:currenteq}
For each \(\sigma\in\{\Car,\Cur,\Hom\}\), Definition \ref{def:fixed} and Definition \ref{def:deep} imply
\begin{equation}
\PrimCurrent{\sigma}\circ\ChainSelect{\sigma}
=
\CurrentExtract{\sigma}\circ\DeepCurrentCarrier{\sigma},
\label{eq:deepcurrent}
\end{equation}
\begin{equation}
\PrimTheta{\sigma}\circ\ChainSelect{\sigma}
=
\EqExtract{\sigma}\circ\DeepEqCarrier{\sigma},
\label{eq:deepeq}
\end{equation}
and therefore
\begin{equation}
(\mathbf 1-\EqProj{\sigma})\PrimTheta{\sigma}\circ\ChainSelect{\sigma}
=
(\mathbf 1-\EqProj{\sigma})\EqExtract{\sigma}\circ\DeepEqCarrier{\sigma}.
\label{eq:deepproj}
\end{equation}
\end{proposition}

\begin{proof}
Using \eqref{eq:carriercomposites} and \eqref{eq:carriercompat},
\[
\PrimCurrent{\sigma}\circ\ChainSelect{\sigma}
=
\CurrentExtract{\sigma}\circ\CurrentCarrier{\sigma}\circ\ChainSelect{\sigma}
=
\CurrentExtract{\sigma}\circ\DeepCurrentCarrier{\sigma},
\]
which is \eqref{eq:deepcurrent}. The same argument gives \eqref{eq:deepeq}:
\[
\PrimTheta{\sigma}\circ\ChainSelect{\sigma}
=
\EqExtract{\sigma}\circ\EqCarrier{\sigma}\circ\ChainSelect{\sigma}
=
\EqExtract{\sigma}\circ\DeepEqCarrier{\sigma}.
\]
Applying \((\mathbf 1-\EqProj{\sigma})\) to \eqref{eq:deepeq} yields \eqref{eq:deepproj}.
\end{proof}

\begin{definition}[Deeper induced defect tuple]
\label{def:deepdefect}
Define the deeper defect tuple
\[
\DeepDefectTuple:\DeepShell\to\DefectSpace
\]
by
\begin{equation}
\begin{aligned}
\DeepDefectTuple(u_{\Car},u_{\Cur},u_{\Hom})
\equiv\;
\Bigl(
&\CurrentExtract{\Car}\bigl(\DeepCurrentCarrier{\Car}(u_{\Car})\bigr),
\CurrentExtract{\Cur}\bigl(\DeepCurrentCarrier{\Cur}(u_{\Cur})\bigr),
\CurrentExtract{\Hom}\bigl(\DeepCurrentCarrier{\Hom}(u_{\Hom})\bigr),
\\
&(\mathbf 1-\EqProj{\Car})\EqExtract{\Car}\bigl(\DeepEqCarrier{\Car}(u_{\Car})\bigr),
(\mathbf 1-\EqProj{\Cur})\EqExtract{\Cur}\bigl(\DeepEqCarrier{\Cur}(u_{\Cur})\bigr),
\\
&
(\mathbf 1-\EqProj{\Hom})\EqExtract{\Hom}\bigl(\DeepEqCarrier{\Hom}(u_{\Hom})\bigr)
\Bigr).
\end{aligned}
\label{eq:deepdef}
\end{equation}
\end{definition}

\begin{theorem}[The deeper defect tuple is the pullback of the primitive defect tuple]
\label{thm:defectpullback}
The tuple of Definition \ref{def:deepdefect} satisfies
\begin{equation}
\DeepDefectTuple
=
\PrimDefectTuple\circ\ChainSelectTriple.
\label{eq:defectpullback}
\end{equation}
\end{theorem}

\begin{proof}
Let \(\mathbf u=(u_{\Car},u_{\Cur},u_{\Hom})\in\DeepShell\). Using \eqref{eq:deepcurrent} and \eqref{eq:deepproj} sectorwise inside \eqref{eq:primtuple}, we obtain
\[
\PrimDefectTuple(\ChainSelectTriple(\mathbf u))
=
\DeepDefectTuple(\mathbf u).
\]
This is exactly \eqref{eq:defectpullback}.
\end{proof}

\section{Observer-Sector Outputs From the Deeper Defect Tuple}

\begin{definition}[Deeper observer-sector outputs]
\label{def:deepoutputs}
Define the deeper observer-sector outputs by pullback along the total selector:
\[
\DeepResidualSector
\equiv
\PrimResidualSector\circ\ChainSelectTriple,
\qquad
\DeepCompletionSector
\equiv
\PrimCompletionSector\circ\ChainSelectTriple,
\qquad
\DeepMatterSector
\equiv
\PrimMatterSector\circ\ChainSelectTriple.
\]
\end{definition}

\begin{theorem}[The deeper observer-sector outputs factor through the deeper defect tuple]
\label{thm:deepsector}
The outputs of Definition \ref{def:deepoutputs} satisfy
\begin{equation}
\DeepResidualSector
=
\ResidualRead\circ\DeepDefectTuple,
\label{eq:deepres}
\end{equation}
\begin{equation}
\DeepCompletionSector
=
\CompletionRead\circ\DeepDefectTuple,
\label{eq:deepcmp}
\end{equation}
and
\begin{equation}
\DeepMatterSector
=
\MatterRead\circ\DeepDefectTuple.
\label{eq:deepmat}
\end{equation}
Hence the defect-fiber factorization of the observer sector is also induced from the deeper substrate package.
\end{theorem}

\begin{proof}
By \eqref{eq:sectorfactor} and Theorem \ref{thm:defectpullback},
\[
\DeepResidualSector
=
\PrimResidualSector\circ\ChainSelectTriple
=
\ResidualRead\circ\PrimDefectTuple\circ\ChainSelectTriple
=
\ResidualRead\circ\DeepDefectTuple,
\]
which proves \eqref{eq:deepres}. The proofs of \eqref{eq:deepcmp} and \eqref{eq:deepmat} are identical.
\end{proof}

\section{One Operative Metric and Ordinary Matter From the Deeper Package}

\begin{definition}[Deeper metric and ordinary matter outputs]
\label{def:deepmetric}
Define
\[
\DeepMetricOut:\DeepShell\to\{\text{observer metrics}\}
\]
and
\[
\DeepMatterOut:\DeepShell\times\Psi\to\{\text{observer stress tensors}\}
\]
by
\begin{equation}
\DeepMetricOut
\equiv
\MetricOut\circ\ChainSelectTriple
\label{eq:deepmetricdef}
\end{equation}
and
\begin{equation}
\DeepMatterOut(\mathbf u,\psi)
\equiv
\MatterOut(\ChainSelectTriple(\mathbf u),\psi).
\label{eq:deepmatterdef}
\end{equation}
\end{definition}

\begin{theorem}[The deeper package yields the same one operative metric and ordinary matter route]
\label{thm:deepmetric}
For every \(\mathbf u\in\DeepShell\) and \(\psi\in\Psi\),
\begin{equation}
\DeepMetricOut(\mathbf u)=\CompletedMetric
\label{eq:deepmetric}
\end{equation}
and
\begin{equation}
\DeepMatterOut(\mathbf u,\psi)
=
T_{\mu\nu}^{\mathrm{matter}}[\CompletedMetric,\psi].
\label{eq:deepmatter}
\end{equation}
Moreover, \(\DeepMetricOut\) and \(\DeepMatterOut\) are the unique deeper outputs induced by the chain-faithful selector, and the deeper package introduces neither a second operative metric nor an enlarged ordinary matter law.
\end{theorem}

\begin{proof}
Equation \eqref{eq:deepmetric} follows immediately from \eqref{eq:deepmetricdef} and \eqref{eq:fixedoutputs}:
\[
\DeepMetricOut(\mathbf u)
=
\MetricOut(\ChainSelectTriple(\mathbf u))
=
\CompletedMetric.
\]
Likewise,
\[
\DeepMatterOut(\mathbf u,\psi)
=
\MatterOut(\ChainSelectTriple(\mathbf u),\psi)
=
T_{\mu\nu}^{\mathrm{matter}}[\CompletedMetric,\psi],
\]
which is \eqref{eq:deepmatter}. The uniqueness statement is immediate from Definitions \ref{def:deepmetric} and \ref{def:fixed}: once the total selector is fixed, the induced deeper outputs are exactly the pullbacks of the already unique primitive outputs. Since the primitive outputs carry no second metric and no enlarged ordinary matter law, neither do their pullbacks.
\end{proof}

\section{Main Theorem}

\begin{theorem}[Primitive-reservoir observer-reduction chain forcing theorem]
\label{thm:main}
Assume the fixed continuation data of Definition \ref{def:fixed} and a chain-faithful deeper substrate package in the sense of Definition \ref{def:deep}. Then the current primitive observer-reduction chain is forced by deeper explicit substrate structure rather than separately fixed. More precisely:
\begin{enumerate}
    \item the sectorwise forcing selectors and the total selector are uniquely determined by the deeper visible maps and deeper carrier data;
    \item the current / projected-equilibrium constituent of the primitive observer-reduction datum is induced from the deeper current and equilibrium carriers by the already fixed shell extractors;
    \item the observer-sector outputs on the deeper package factor through the deeper induced defect tuple;
    \item the deeper metric output is exactly the same one operative metric \(\CompletedMetric\);
    \item the deeper ordinary matter output is exactly the same standard matter route through \(\CompletedMetric\);
    \item consequently the deeper forcing introduces neither a second operative metric nor an enlarged low-energy matter law.
\end{enumerate}
Therefore the positive gain is a genuine deeper derivational step below the current primitive observer-reduction datum: the existence of the now-recorded chain is justified at current scope by deeper explicit substrate structure, while the benchmark one-metric claim boundary remains fixed.
\end{theorem}

\begin{proof}
Item (1) is Proposition \ref{prop:selector_unique} together with Corollary \ref{cor:selector_total}. Item (2) is Proposition \ref{prop:currenteq}. Item (3) is Theorem \ref{thm:deepsector}. Items (4) and (5) are Theorem \ref{thm:deepmetric}. Item (6) is the uniqueness content of Theorem \ref{thm:deepmetric}. The final sentence is the combined routing consequence: this note no longer treats the primitive observer-reduction chain as fixed datum, but it also does not change the benchmark one-metric package that the chain produces.
\end{proof}

\begin{corollary}[Why the present gain is not another same-output relay]
\label{cor:notrelay}
The positive content of Theorem \ref{thm:main} is not merely that the same operative metric and the same ordinary matter route reappear. Its benchmark-facing gain is that the chain producing those outputs has itself moved from recorded datum to deeper forced output.
\end{corollary}

\begin{proof}
By Theorem \ref{thm:main}, the same one-metric and ordinary-matter package is preserved, but the selector, current / projected-equilibrium package, defect-fiber factorization, and output maps are now all induced by the deeper substrate package. The novelty therefore lies in the derivational status of the chain itself, not in a change of low-energy outputs.
\end{proof}

\begin{corollary}[What remains open after this note]
\label{cor:next}
After Theorem \ref{thm:main}, the remaining honest burden is no longer to defend the current primitive observer-reduction chain as an acceptable fixed datum. The remaining burden is to realize or derive the chain-faithful deeper substrate package itself from still more explicit substrate equations, selection principles, or stationary structure, while preserving the same benchmark one-metric package.
\end{corollary}

\begin{proof}
Theorem \ref{thm:main} already shows that once a chain-faithful deeper substrate package is given, the primitive observer-reduction chain is forced and the same one-metric package follows without enlargement. What remains open is therefore the deeper origin of that forcing package itself, not another same-layer note about the outputs it already induces.
\end{proof}

\section{Conclusion}

The positive result of this note is narrow but real. The current primitive observer-reduction chain is no longer treated here as merely recorded datum once a chain-faithful deeper substrate package is supplied. Under explicit visible-map and carrier compatibility together with the chain-separation hypothesis, the sectorwise forcing selectors are unique, the current / projected-equilibrium package is induced from deeper carriers, the observer-sector outputs factor through the deeper defect tuple, and the deeper metric and ordinary matter outputs are exactly the same one operative metric \(\CompletedMetric\) and the same standard matter route through that metric.

The scope remains disciplined. This note does not claim a completed first-principles derivation of GR from final substrate microphysics, and it does not yet derive the deeper forcing package from a final substrate action. It also does not enlarge the benchmark observable sector beyond the exact one-metric package already in hand. What it does do is move the derivational burden one honest layer lower: the present primitive observer-reduction chain is now justified as a forced consequence of deeper explicit substrate structure rather than simply carried as fixed datum.

\end{document}
